\subsection{课题目标}
\label{sec:purpose}

本课题研究低时延大语言模型推理中相邻算子之间的执行方式。计算图上的每一条边都面临同一个选择：生产者与消费者留在同一 grid 内，还是分属两个 grid、在二者之间形成物理边界；若形成物理边界，后继是否借助 Programmatic Dependent Launch（PDL）提前启动，与前驱重叠。开题报告把这一选择表述为全图执行计划的优化
\begin{equation}
\Pi^{\star}(G,x)=\arg\min_{\Pi\in\mathcal{F}(G,x)} T(G,x;\Pi), \label{eq:obj}
\end{equation}
其中 $G=(O,E)$ 为计算图，$x$ 为输入规模与运行配置，$\Pi$ 为各边处理方案 $\pi(e)=\langle m_e,\theta_e\rangle$ 的集合，$\mathcal{F}(G,x)$ 为满足数据依赖、内存可见性与硬件资源约束的合法计划集合。围绕式~(\ref{eq:obj})，开题报告确定了四项研究内容：计算图依赖关系建模、依赖感知的 PDL 跨边界重叠、全图边处理方案搜索，以及框架实现与端到端验证。

\subsection{计划执行情况}

按开题计划，至 2026 年 11 月应完成前两项研究内容并启动第三项。实际执行情况如表~\ref{tab:plan} 所示：前两项的主体工作已按计划完成；第三项完成了 decoder layer 链上的规划器及其与固定方案的实测比较；第四项提前启动，已建立统一算子后端，并贯通 prefill 与单步 decode 的正确性通路。

\begin{table}[htbp]\centering
\caption{研究计划执行情况}\label{tab:plan}
\small
\begin{tabularx}{\textwidth}{@{}p{3.3cm}p{2.1cm}X@{}}
\toprule
研究内容 & 计划时间 & 完成情况\\
\midrule
计算图依赖关系建模 & 2026.07–08 & 基本完成。建立阶段级依赖抽象、逐边合法候选与 grid 资源包络模型；task/tile 级就绪粒度在 Hazy Megakernels 原型上完成验证；面向任意计算图的 lowering 尚未形成独立模块。\\
依赖感知的 PDL 跨边界重叠 & 2026.09–10 & 基本完成。实现触发、同步与权重预取的参数化 PDL 交接，完成逐边 PDL 选择与边级实测，给出后继提前进入的条件、落位的派发模型与不依赖逐边实测的选择判据；触发时机与预取量目前取固定值。\\
全图边处理方案搜索 & 2026.11–12 & 进行中。完成 32 种分组方式的枚举式规划器与实测校验；建立矩阵阶段、attention 与多阶段分组的时延模型并接入规划器，以约四成的测量量得到时延相当的方案；在 15 个负载上完成各项决策对输入规模依赖的分类；层间异构方案、候选计划表与在线选择尚未完成。\\
框架实现与端到端验证 & 2027.01–04 & 提前启动。完成统一算子后端与 prefill、单步 decode 的数值校验；连续 token 生成与 TPOT 分位数尚未完成。\\
\bottomrule
\end{tabularx}
\end{table}

相对开题报告，研究过程中作了两点调整。其一，边处理参数 $\theta_e$ 由同步与 PDL 参数扩展到 grid 级资源配置。实验表明，一条边是否应成为物理边界，主要取决于两侧算子共享同一 launch 资源包络所付出的代价（第~\ref{sec:hw}~节），因此线程块数、分块形状与分阶段驻留配置须与边界选择一并决定。其二，评测目标由 decode 单步扩展到 prefill 与 decode 两个阶段。两个阶段中各算子的时延占比与资源需求不同，同一条边的最优处理方式随之变化，仅以稳态 decode 评测不足以反映首 token 时延。

\subsection{已取得的阶段性成果}

\begin{enumerate}[leftmargin=2.6em,itemsep=2pt]
\item 建立了阶段级的边处理抽象与 grid 资源包络模型，将分块形状、线程块数与分阶段驻留配置纳入边处理参数，并以共驻代价给出是否切分的判据。
\item 完成了 H100 上边界交接代价、访存带宽、阶段并行度需求、线程块放置与主循环构成的定量分析，查明小分块的主循环受访存一侧限制。
\item 实现了统一算子后端、同 grid 内的分阶段驻留配置与参数化的 PDL 交接，全部方案通过数值校验。
\item 实现了逐边规划器。相对充分调优的单一 grid，规划方案在 decode 上最高降低 13.0\% 的单步时延，在 prefill 上最高降低 19.2\% 的首 token 时延；在资源需求相近的负载上规划器保留单一 grid。
\item 建立了矩阵阶段、attention 与多阶段分组的时延模型：矩阵阶段在 40 个留出行数上的中位误差为 3.4\%–3.9\%，分组模型在 2177 组实测配置上为 4.9\%。
\item 给出了 PDL 下以 carve-out 档位与资源余量表述的准入条件、档位沿发射链保持的规律与后继落位的派发模型，并据此划分三类边界；后继以矩阵阶段开头时，判为均衡的 113 类实测边界中 94\% 因 PDL 加速。加入 L1 判据后，按判据选择 PDL 与逐边搜索所得的时延降幅中位数相近。
\item 把代价模型与 PDL 判据接入规划器，以穷举式规划 43\% 的测量量得到时延相差不超过 1.0\% 的方案。
\item 在 15 个负载、1649 个端到端方案上测量了各项决策对输入规模的依赖，据此把决策分为固定型、切换型与参数型：decode 上边的切分可以离线固定，prefill 上只需两个机制候选（最大损失 0.6\%），各组的资源配置须按约 22 个行数区间确定；并给出离线准备、在线查表的系统结构。
\end{enumerate}


\subsection{存在的问题}\label{sec:issues}

\begin{enumerate}[leftmargin=2.6em,itemsep=2pt]
\item \textbf{收益集中于部分负载。}batch 为 1 的短上下文 decode 与最短的 prefill 中，各阶段的资源需求相近，单一 grid 仍是最优方案，规划方案的收益不超过 0.3\%。这些负载恰是低时延场景的典型工作点；其上可能的改进来自同 grid 内的细粒度事件范式，而该范式尚未并入统一后端。
\item \textbf{算子实现仍有差距。}batch 为 1 的 decode 比 Hazy Megakernels 慢 26\%；行数不小于 256 的 prefill 中，后端的 GEMM 落后于 cuBLAS，规划方案比逐算子 CUDA Graph 慢 28\%–73\%。边处理方式的结论建立在同一后端内部的比较之上，但端到端的绝对时延受此限制。
\item \textbf{依赖抽象的粒度与通用性不足。}当前抽象停留在阶段级，分组方式各层相同，候选空间为 32 种；task/tile 级的就绪事件只在 Hazy Megakernels 原型上验证过，面向任意计算图的 lowering 尚未形成独立模块。
\item \textbf{模型的精度只够排序。}分组模型在含 attention 的 prefill 分组上的中位误差为 8.0\%–15.3\%，对分阶段驻留配置的估计偏乐观；模型引导的规划器因此仍需实测排在前列的配置与全部单阶段分组，测量量为穷举式的 43\%，个别负载上所得方案比穷举式方案慢 1.0\%。含 embedding 与词表投影的分组尚无模型，这些分组仍测全部配置。
\item \textbf{集中放置的减速幅度尚无解析式。}派发模型可以预测后继的落位，探针作为 PDL 后继运行时也能复现真实方案中 U 的减速（5.3–7.2 对 7.6\us），但减速随各 SM 上线程块数的分布而变，目前只有“读取带宽在 SM 之间的分配不与线程块数成正比”这一定性规律，没有由分布直接给出时延的公式；真实方案中部分层的后继均匀地落在 64 个 SM 上（每个 SM 2 个），执行时间仍增加约三成，这一部分尚未解释。规划时这不影响取舍，因为可能被集中放置的边界不启用 PDL。可共驻的边界上，运行时常使少数 SM 多驻留一个线程块，模型对这类边界的命中率为 134/172。
\item \textbf{L1 敏感度按负载取得。}attention 对 L1 容量的敏感度随序列长度变化，目前由每个负载上一次强制档位的运行测得，尚未写成序列长度的函数；carve-out 作为分组的配置项也尚未进入搜索。attention 分组每个 SM 只能驻留 4 个线程块的原因仍未查明。
\item \textbf{硬件性能计数器不可用。}实验平台的计数器仅对管理员开放，主循环内逐指令的停顿归因无法取得，图~\ref{fig:ablate} 的消融只能给出部件一级的结论。
\item \textbf{评测负载的行数未覆盖分块边界。}第~\ref{sec:exp}~节的 15 个负载的行数均为分块行数的整倍数，而行数刚越过整倍数时矩阵阶段的时延可升高 47\%（第~\ref{sec:cm}~节），真实请求的行数并不受此约束。
\item \textbf{端到端验证不完整。}评测限于 prefill 与单步 decode，尚未接入连续 token 生成，也未报告 TPOT 与 P50/P90/P99；模型与硬件各只有一种。
\item \textbf{决策分类的覆盖范围有限。}三类决策的划分建立在 15 个负载之上，矩阵阶段的行数区间取自合成驱动上的单阶段分组；二者尚未扩展到 batch 与序列长度的网格和多阶段分组。候选计划表与在线阶段尚未实现，查表与切换计划的开销没有实测数据。
\end{enumerate}

